\documentclass{article}
\begin{document}

\section{Exercício 1 - Isolamento de raízes}

\subsection{Contexto - Teorema de Bolzano}
    \subsubsection{Explicação: Analogia da Ponte}
        Pense no teorema como uma ponte sobre um rio: se você começa em uma margem (onde $f(a)$ tem um sinal) e termina na outra margem (onde $f(b)$ tem o sinal oposto), e a ponte é contínua (a função é contínua), você obrigatoriamente cruzou o rio (passou pelo ponto onde $f(x) = 0$)

        Ou seja, o teorame garante a existência de pelo menos uma raiz, podendo haver um número ímpar de raízes. Porém, se $f(a)$ e $f(b)$ têm o mesmo sinal, pode haver um número par de raízes ou nenhuma, e não é possível usar esse teorema.

    \subsubsection{Pass-a-passo}
        \begin{enumerate}
            \item \textbf{Verifique a Continuidade:} Certifique-se de que a função $f(x)$ é contí\-nua no intervalo que você está analisando. A maioria das funções polinomiais, exponenciais, logarítmicas e trigonométricas são contínuas em seus domínios.
            \item \textbf{Escolha um Intervalo:} Selecione um intervalo $[a,b]$ onde você suspeita que possa haver uma raiz (a análise gráfica pode ser um ótimo ponto de partida aqui!)
            \item \textbf{Calcule os Valores nos Extremos:} Avalie a função nos pontos $a$ e $b$, ou seja, calcule $f(a)$ e $f(b)$.
            \item \textbf{Verifique a Mudança de Sinal:} Observe os sinais de $f(a)$ e $f(b)$. Se eles forem opostos (um positivo e outro negativo), então o Teorema de Bolzano garante que existe pelo menos uma raiz no intervalo $(a,b)$.
        \end{enumerate}  

\subsection{a) $f(x) = x^3 - 2x - 5$}

    \subsubsection{Passo 1 - Verificar a Continuidade}
        A função $f(x) = x^3 - 2x - 5$ é um polinômio, portanto é contínua em todo os Reais.
    \subsubsection{Passo 2 - Escolha um Intervalo}
            Intervalo aleatório escolhido: $[-5, 5]$

    \subsubsection{Passo 3 - Calcule os Valores nos Pontos Extremos}
        \[
            f(-5) = (-5)^3 - 2*(-5) - 5 = - 125 + 10 - 5 = - 115 - 5 = - 120
        \]
        \[
            f(5) = (5)^3 - 2*(5) - 5 = 125 - 10 - 5 = 115 - 5 = 110
        \]

\subsubsection{Passo 4 - Verificar a Mudança de Sinal}
        $f(-5) < 0$ e $f(5) > 0$, logo existe $x \in [-5,5]$ tal que $f(x) = 0$
\section{Exercício 2 - Método da Bissecção}

\section{Exercício 3 - Método da Posição Falsa}

\section{Exercício 1 - Método das Aproximações Sucessivas}

\section{Exercício 1 - Método de Newton}

\section{Exercício 1 - Método da Secante}

\section{Exercício 1 - Comparação de Métodos}


\end{document}